\section{Triangle Inequality} \label{sec: tri_ineq}

Below is the proof that the adjusted correlation metric $f$ satisfies an approximate triangle inequality. 

\begin{proof}[Proof of Lemma \ref{lem: adjusted_tri_ineq}]
    Note that after Step 1 of Definition \ref{def: adj_corr_metric}, $f$ satisfies the triangle inequality, since $f=d$ here and for any $u,v,w$, $d_{uv} \leq d_{uw}+d_{vw}.$
    We then consider the execution of Step 2 in Definition \ref{def: adj_corr_metric}.
    The triangle inequality on the adjusted correlation metric holds if either of $(u,w)$ or $(w,v)$ have distance 1 with respect to the adjusted correlation metric, and it also clearly holds if the three edges have the same distance as in the correlation metric. 
    It follows that the only case to check is when $f_{uv}=1$, $f_{uw}=d_{uw}$, and $f_{wv}=d_{wv}$.
    If $f_{uv}=1$ at the end of Step 2, by the definition of $f$ it must be that $d_{uv} > 0.7$.
    Since $d_{uv} \leq d_{uw}+d_{vw}$, we have that $0.7 \leq d_{uw}+d_{vw}$, 
    and so $f_{uv}=1 \leq \frac{10}{7}(d_{uw}+d_{vw}) \leq \frac{10}{7}(f_{uw}+f_{vw})$ at the end of Step 2.
    
    Finally, Step 3 of Definition \ref{def: adj_corr_metric} also preserves this approximate triangle inequality. 
    The $f$ values on the edges leaving a vertex are raised to 1 together, and any triangle involving one of these edges necessarily involves another.
\end{proof}

\section{$\ell_{\infty}$-norm Approximation} \label{sec: infty_norm}
Below we include the proof that the adjusted correlation metric $f$ has bounded fractional cost in the $\ell_{\infty}$-norm. 

\begin{proof}[Proof of Lemma \ref{lem: infty_proof}]
    The fractional cost of the correlation metric $d$ is at most $8 \cdot \textsf{OPT}$ by \cite{DMN23}. That is,
    \begin{equation} \label{eq: frac_cost_infty}
        \sum_{v \in N_u^+} d_{uv} + \sum_{v \in N_u^-} (1-d_{uv}) \leq 8 \cdot \textsf{OPT} \hspace{0.3cm} \text{for all } u \in V.
    \end{equation}

    Consider any $u \in V$. Observe that the second sum in (\ref{eq: frac_cost_infty}) can only decrease if Step 2 in Definition \ref{def: adj_corr_metric} is executed for any of the negative edges incident to $u$. Now let $R_1$ denote the set of vertices $u$ such that Step 3 is executed on $u$. For $u \in R_1$, the  second sum decreases by at least $\frac{3}{10} \cdot \frac{10}{3} \cdot \Delta_u = \Delta_u $, since there are at least $\frac{10}{3}\cdot \Delta_u$ negative edges $(u,v)$ with $1-d_{uv} \geq \frac{3}{10}$, and these edges are all rounded up  to 1 in $f$, so $1-d_{uv}$ is rounded down to 0. On the other hand, the first sum increases by at most $\Delta_u$, since there are $\Delta_u$ positive edges incident to $u$ and each gets rounded up to 1 in $f$. Therefore, we know that (\ref{eq: frac_cost_infty}) does not increase when we replace $d$ with $f$ when $u \in R_1$. 

Finally, it is possible that (\ref{eq: frac_cost_infty}) increases even when $u \not \in R_1$. This is because $u$ may have neighbors $v$ such that $v \in R_1$. For such $v$, $d_{uv}$ is raised to $f_{uv}=1$. This is only better for the second sum in (\ref{eq: frac_cost_infty}). Now we consider what happens to the first sum in (\ref{eq: frac_cost_infty}) when we replace $d$ with $f$. Let $R_2 = V \setminus R_1$. Fix $u \in R_2$. Then
\begin{align*}
    \sum_{v \in N_u^+} f_{uv} &= \sum_{v \in N_u^+ \cap R_1} 1 + \sum_{v \in N_u^+ \cap R_2} d_{uv} \\
    &= \sum_{\substack{v \in N_u^+ \cap R_1 \\ d_{uv} \leq 1/4}} 1 + \sum_{\substack{v \in N_u^+ \cap R_1 \\ d_{uv} > 1/4}} 1 + \sum_{v \in N_u^+ \cap R_2} d_{uv} \\
    &\leq \sum_{\substack{v \in N_u^+ \cap R_1 \\ d_{uv} \leq 1/4}} 1 + \sum_{\substack{v \in N_u^+ \cap R_1 \\ d_{uv} > 1/4}} 4 \cdot d_{uv} + \sum_{v \in N_u^+ \cap R_2} d_{uv} \\
    &\leq \sum_{\substack{v \in N_u^+ \cap R_1 \\ d_{uv} \leq 1/4}} 1 + 4 \cdot \sum_{v \in N_u^+} d_{uv}
\end{align*}
The second term above is bounded by $32 \cdot \textsf{OPT}$ by (\ref{eq: frac_cost_infty}). To bound the first term, let $w^*$ be an arbitrary vertex in $N_u^+ \cap R_1$ (if this set is empty, we just ignore the first term). Since $d_{uw^*} \leq 1/4$, we know by Proposition \ref{prop: similar_nbhds} that $\Delta_u \leq \frac{7}{3} \cdot \Delta_{w^*}$. Also, 
$$ \sum_{v \in N_{w^*}^+} f_{w^*v} + \sum_{v \in N_{w^*}^-} (1-f_{w^*v})  \leq \Delta_w^*$$ 
since $w^* \in R_1$, and by the argument in the previous paragraph we have that $\Delta_{w^*} \leq 8 \cdot \textsf{OPT}$. So
\[\sum_{\substack{v \in N_u^+ \cap R_1 \\ d_{uv} \leq 1/4}} 1 \leq \Delta_u \leq \frac{7}{3} \cdot \Delta_{w^*} \leq \frac{7}{3} \cdot 8 \cdot \textsf{OPT}.\] This concludes the proof. 
\end{proof}
  

\section{$\ell_1$-norm Approximation}\label{sec:l1-norm}
The next lemma proves that the adjusted correlation metric well approximates the optimal in the $\ell_1$-norm.
We will show that the fractional cost of the adjusted correlation metric $f$ in the $\ell_1$-norm is a constant factor approximation to the optimal, i.e., 
    $$\sum_{u \in V} \sum_{v \in N_u^+} f_{uv} + \sum_{u \in V} \sum_{v \in N_u^-} (1-f_{uv}) = O(\textsf{OPT}).$$
We argue the second sum---over the negative edges---has a constant factor approximation in the first lemma (Lemma \ref{lem: l1_bdd-neg}), then show the first sum---over the positive edges---does too in the second lemma (Lemma \ref{lem: l1_bdd-pos}). 
Combining these lemmas gives the following approximation guarantee for the $\ell_1$-norm.


\begin{lemma}
    The fractional cost of the adjusted correlation metric $f$ in the $\ell_1$-norm objective is a constant factor away from the cost of the optimal integral solution in the $\ell_1$-norm
\end{lemma}
\begin{proof}
The fractional cost in the $\ell_1$-norm is
     $$\sum_{u \in V} \sum_{v \in N_u^+} f_{uv} + \sum_{u \in V} \sum_{v \in N_u^-} (1-f_{uv}) \leq 74 \cdot \textsf{OPT}.$$
The first sum is bounded by $34\cdot \textsf{OPT}$ by Lemma \ref{lem: l1_bdd-pos} and the second sum is bounded by $40\cdot \textsf{OPT}$ by Lemma \ref{lem: l1_bdd-neg}. 
\end{proof}



\begin{lemma} \label{lem: l1_bdd-pos}
    The fractional cost of the adjusted correlation metric $f$ in the $\ell_1$-norm objective for the set of positive edges is a constant factor approximation to the optimal, i.e.,
    $$ \sum_{u \in V} \sum_{v \in N_u^+} f_{uv} = 34 \cdot \textsf{OPT}.$$
\end{lemma}

\begin{proof}[Proof of Lemma \ref{lem: l1_bdd-pos}]


Fix an optimal clustering $\mathcal{C}$. 
We partition vertices based on membership in $C(u)$ or $\overline{C(u)}$ and let $y$ denote the disagreement vector of $\mathcal{C}$. 
We prove in 
the following claim that the total fractional cost of edges $e \in E^+$ such that $f_e = d_e$ is bounded.  
Then, it just remains to bound the fractional cost of positive edges whose distances are raised to 1 in Step 3 of Definition \ref{def: adj_corr_metric}.
\begin{claim}\label{claim: l1-pos-cm}
The total fractional cost of positive edges for the $\ell_1$-norm objective when
$f_e = d_e$ is a $3$-approximation. 
\end{claim}
\begin{proof}[Proof of Claim \ref{claim: l1-pos-cm}]
For the $\ell_1$-norm, $\textsf{OPT} = \sum_{u \in V} y(u)$.
So for a fixed vertex $u$,
we partition the fractional cost of disagreements from positive edges incident to $u$, i.e. $(u,v) \in E^+$, based on whether vertices $v$ are in $C(u)$ or not:
\begin{align*}
    \sum_{u \in V} \sum_{v \in N_u^+} f_{uv}
    &=  \sum_{u \in V}\sum_{v \in N_u^+ \cap C(u)} d_{uv}+ \sum_{u \in V}\sum_{v \in N_u^+ \cap \overline{C(u)}} d_{uv} \\
    &= \sum_{u \in V}\sum_{v \in N_u^+ \cap C(u)} \frac{|N_u^+ \cap N_v^-| + |N_u^- \cap N_v^+|}{n-|N_u^- \cap N_v^-|} + \sum_{u \in V}\sum_{v \in N_u^+ \cap \overline{C(u)}} d_{uv}. \notag 
\end{align*}
The right-hand sum is easy as we can just use that $d_{uv}$ is at most 1. 
$$\boxed{\sum_{u \in V}\sum_{v \in N_u^+ \cap \overline{C(u)}} d_{uv} \leq \sum_{u \in V}\sum_{v \in N_u^+ \cap \overline{C(u)}} 1 \leq \sum_{u \in V} y(u) = \textsf{OPT}}
$$

For the left-hand sum, we will exploit that $d_{uv}$ is symmetric, and specifically that we can bound the denominator by $|N_u^+|$ or $|N_v^+|$.  
\begin{align*}
    &\sum_{u \in V}\sum_{v \in N_u^+ \cap C(u)} \frac{|N_u^+ \cap N_v^-| + |N_u^- \cap N_v^+|}{n-|N_u^- \cap N_v^-|}
    \leq \sum_{u \in V} \sum_{v \in N_u^+ \cap C(u)}\frac{y(u) + y(v)}{n-|N_u^- \cap N_v^-|}\\
    &\leq \sum_{u \in V} \sum_{v \in N_u^+} \frac{y(u) + y(v)}{n-|N_u^- \cap N_v^-|} 
    = \sum_{u \in V} \sum_{v \in N_u^+} \frac{y(u)}{n-|N_u^- \cap N_v^-|} + \sum_{u \in V} \sum_{v \in N_u^+} \frac{y(v)}{n-|N_u^- \cap N_v^-|}\\
    &= \sum_{u \in V} \sum_{v \in N_u^+} \frac{y(u)}{n-|N_u^- \cap N_v^-|} + \sum_{v \in V} \sum_{u \in N_v^+} \frac{y(v)}{n-|N_u^- \cap N_v^-|} \\
    &\leq \sum_{u \in V} \frac{y(u)}{|N_u^+|} \cdot |N_u^+| + \sum_{v \in V} \frac{y(v)}{|N_v^+|} \cdot |N_v^+| \leq 2\cdot \textsf{OPT}.
\end{align*}
All in all:
$$
    \boxed{\sum_{u \in V}\sum_{v \in N_u^+ \cap C(u)} \frac{|N_u^+ \cap N_v^-| + |N_u^- \cap N_v^+|}{n-|N_u^- \cap N_v^-|}
    \leq 2\cdot \textsf{OPT}}
$$
\end{proof}


Now, we must bound the cost of the positive edges whose distances were raised in Step 3 of Definition \ref{def: adj_corr_metric}. 
Let $R_1 = \{u : \vert\{N_u^- \cap \{v: d_{uv} \leq 0.7\}\vert \geq \frac{10}{3}\cdot \Delta_u\}$.
For any vertex $u$, define the shorthand $R_1(u) = \{N_u^- \cap \{v: d_{uv} \leq 0.7\}\}$
\footnote{The use of a subscript may not be clear here, but it is helpful in the $\ell_p$-norm proof.}; 
recall that $u \in R_1$ is put in its own cluster by the adjusted correlation metric.
Then, our next goal is to show that
\begin{equation} \label{eq: ell1_R1}
2 \cdot \sum_{u \in R_1} |N_u^+| = O(\textsf{OPT}).
\end{equation}
Note that the factor of 2 comes from the observation that even if $u \not \in R_1$, $u$ may have neighbors $v \in R_1 \cap N_u^+$, in which cases $f_{uv}$ is raised to 1 even though Step 3 is not executed on $u$. Fix a $u \in R_1$.
We subdivide into several cases based on $\mathcal{C}$.

\paragraph{Case 1: }\emph{At least a 0.15 fraction of $N_u^+$ is in clusters other than $C(u)$.} \\
Let $ u \in V^{1}$ be the vertices in this case.
There are at least $0.15\cdot |N_u^+|$ disagreements incident to $u$, so 
$$\boxed{\sum_{u \in V^1} |N_u^+| \leq \sum_{u \in V^1} \frac{1}{0.15} \cdot y(u) \leq \frac{20}{3}\cdot \textsf{OPT}.}$$

\paragraph{Case 2: }\emph{At least a 0.85 fraction of $N_u^+$ is in $C(u)$.} \\
We subdivide Case 2 into two cases depending on how much 
of $R_1(u)$ is in $\overline{C(u)}$.\\

\noindent \textbf{Case 2a:} \emph{At least half of $R_1(u)$ is in clusters other than $C(u)$.} \\
Let $ u \in  V^{2a}$ be the vertices in this case.
Denote this subset of $R_1(u)$ by $R'_1(u) = R_1(u) \cap \overline{C(u)}$, and note that $|R'_1(u)| \geq \frac{5}{3}\cdot \Delta_u$. If the number of negative edges incident to $u$ that are contained in $C(u)$ is more than $\Delta_u$, then charge $|N_u^+| = \Delta_u$ to $y(u)$, since these negative edges are disagreements. These $u \in V^{2a}$ thus contribute $\textsf{OPT}$ to (\ref{eq: ell1_R1}). 

Otherwise, $|C(u)| \leq 2\Delta_u$.  Consider $v \in R'_1(u)$. Again, let $N_{u,v} = N_u^+ \cap N_v^+ \cap C(u)$. Since $d_{uv} \leq 0.7$ and by the assumption bringing us into Case 2, we apply Proposition \ref{prop: pos_overlap} to see that $|N_{u,v}| \geq 0.15 \cdot |N_u^+|$. 
For each $v \in R'_1(u)$ and $w \in N_{u,v}$, the positive edge $(w,v)$ is a disagreement. 


   For $u \in  V^{2a}$ and $w \in N_u^+ \cap C(u)$, define 
    $$\varphi(u,w) = |R_1'(u)\cap N_w^+  |.$$
    Let each $w \in N_u^+ \cap C(u)$ dispense charge $\varphi(u,w)/|C(u)|$ to $u$.
First, observe that $u$ receives $\Omega(|N_u^+|)$ charge; specifically, $u$ receives charge  at least 
 \begin{align*}
    \sum_{w \in N_u^+ \cap C(u)} \frac{\varphi(u,w)}{|C(u)|}
    &= \sum_{w \in N_u^+ \cap C(u)} \frac{|R_1'(u)  \cap N_w^+|}{|C(u)|}  \\
       & = \frac{1}{|C(u)|}\cdot\sum_{w \in N_u^+ \cap C(u)} \sum_{v \in R_1'(u)  \cap N_w^+} 1  
    = \frac{1}{|C(u)|}\cdot\sum_{v \in R_1'(u)} \sum_{\underset{ \cap C(u) \cap N_u^+}     {w \in  \cap N_v^+}} 1 \\
        &= \frac{1}{|C(u)|}\cdot\sum_{v \in R_1'(u)} | C(u) \cap N_u^+ \cap N_v^+| \geq 
        \frac{1}{|C(u)|}\cdot\sum_{v \in R_1'(u)} 0.15\cdot | N_u^+ | \\
    &\geq \frac{1}{|C(u)|}\cdot 0.15 \cdot |N_u^+| \cdot |R_1'(u)| \geq \frac{1}{2\Delta_u}\cdot 0.15 \cdot \Delta_u \cdot \frac{5}{3}\Delta_u = \frac{1}{8} \cdot \Delta_u.
 \end{align*}


Next we upper bound the charge dispensed in total to \emph{all} $u \in V^{2a}$. Note by definition that $\varphi(u,w) \leq y(w)$. Every vertex $w \in V$ dispenses at most
    $$\frac{y(w)}{|C(u)|} = \frac{y(w)}{|C(w)|}$$
    charge to each $u \in C(w) \cap N_w^+$. 
    So each $w$ dispenses a total of charge at most 
    \[\frac{y(w)}{|C(w)|} \cdot |C(w)| = y(w)\]
     to all $u \in V^{2a}$. 
    Putting the lower and upper bounds together:  
    \begin{align*}
        \sum_{w \in V} y(w) &\geq \text{total charge dispensed} \\
        &\geq \sum_{u \in  V^{2a}: |C(u)| \leq 2\Delta_u} \frac{1}{|C(u)|} \sum_{w \in N_u^+ \cap C(u)} \varphi(u,w) \\
        &\geq \sum_{u \in  V^{2a}: |C(u)| \leq 2\Delta_u} \frac{1}{8} \cdot |N_u^+|
    \end{align*}
    so 
    \[\boxed{\sum_{u \in  V^{2a}} |N_u^+| \leq \textsf{OPT} + 8  \cdot \sum_{w \in V} y(w) \leq 9 \cdot \textsf{OPT}.}\]

\noindent \textbf{Case 2b:} \emph{At least half of $R_1(u)$ is in $C(u)$.}\\
Let $ u \in  V^{2b}$ be the vertices in this case.
There are at least $\frac{5}{3}\Delta_u$ negative disagreeing edges incident to $u$ (those edges $(u,w)$ with $w \in R_1(u)$). So we may charge $|N_u^+|$ to $y(u)$:
$$\boxed{\sum_{u \in V^{2b}} |N_u^+| \leq  \textsf{OPT}}$$
This concludes the bounding of the fractional cost of the negative edges, 
and our constant in the approximation factor is $3+ 20/3 + 9 + 1 < 20$ (adding together the boxed upper bounds).
\end{proof}

\begin{lemma} \label{lem: l1_bdd-neg}
    The fractional cost of the adjusted correlation metric $f$ in the $\ell_1$-norm objective for the set of negative edges is a constant factor approximation to the optimal, i.e.,
    $$ \sum_{u \in V} \sum_{v \in N_u^-} (1-f_{uv}) \leq 40\cdot \textsf{OPT}.$$
\end{lemma}

\begin{proof}[Proof of Lemma \ref{lem: l1_bdd-neg}]
First observe that negative edges that have been rounded up to 1 contribute 0 to the fractional cost. So we may focus on the fractional cost of negative edges $e=(u,v)$ with $d_{uv} = f_{uv} \leq 0.7$. \\


Fix an optimal clustering $\mathcal{C}$, whose value is $\textsf{OPT}$.
If $d_{uv} \leq 0.7$ but $u$ and $v$ are in the same cluster in $\mathcal{C}$,
then $(u,v)$ is a disagreement with respect to $\mathcal{C}$,
so we can charge the fractional cost $1-f_e = 1-d_e \leq 1$ to \boxed{$\textsf{OPT}$}. 

Now we need to consider the negative edges whose endpoints are in different clusters in $\mathcal{C}$.
Recall $C(u)$ is the cluster of $\mathcal{C}$ containing vertex $u$. 
We define the set of \emph{cut edges} incident to $u$ to be
$$E_u = \{(u,v) : (u,v) \in E^-, d_{uv} \leq 0.7, C(u) \neq C(v) \}.$$ 

By Step 3 in the definition of the adjusted correlation metric, $|E_u| \leq 10/3\cdot\Delta_u$. Since the fractional cost of each cut edge is $1-f_{uv} = 1-d_{uv} \geq 0.3 = \Omega(1)$, we need to show that the total number of cut edges is bounded by $\textsf{OPT}$, i.e.,
$\sum_{u \in V} |E_u| = O(\textsf{OPT}).$
We let $V_u = \{v: (u,v) \in E_u\} $,
so our goal is equivalent to showing that 
$\sum_{u \in V}|V_u| = O(\textsf{OPT}).$
The definition of $V_u$ is the same as $R'_1(u)$ in the proof of the previous lemma, 
but $R_1'(u)$ was only defined for $u \in R_1$.


We will charge the value of $\sum_{u \in V} |V_u|$ to disagreements in $\mathcal{C}$. We use $y(u)$ to denote the number of disagreements incident to $u$ with respect to $\mathcal{C}$. 

Fix a vertex $u$. We consider several cases based on the clustering $\mathcal{C}$. 

\paragraph{Case 1: }\emph{At least a 0.15 fraction of $N_u^+$ is in clusters other than $C(u)$.}  \\
Let $V^1$ be the vertices of $V$ in this case. 
For any such $u \in V^1$, 
there are at least $0.15\cdot \Delta_u$ disagreements incident to $u$ in $\mathcal{C}$, 
so $0.15\cdot \Delta_u\leq y(u).$ 
Combining this with the fact that $|V_u| \leq \frac{10}{3} \cdot \Delta_u $, 
we see that
$$|V_u| \leq \frac{10}{3}\cdot \Delta_u \leq \frac{200}{9} \cdot y(u).$$
Summing over all $u \in V^1$, we see that
$$\boxed{\sum_{u \in V^1} |V_u| \leq \frac{200}{9} \cdot \sum_{u \in V} y(u) = \frac{200}{9} \cdot \textsf{OPT}.}$$

\paragraph{Case 2: }\emph{At least a 0.85 fraction of $N_u^+$ is in $C(u)$.} \\
Let $V^2$ be the vertices of $V$ in this case. 
For any such $u \in V^2$, we will show that for every $v \in V_u$, 
there is large overlap between $C(u)$, $N_u^+$, and $N_v^+$. 
Proposition \ref{prop: pos_overlap} will help us as we further refine Case 2 into two settings: 
\medskip

\noindent \textbf{Case 2a:} \emph{At least 0.85 fraction of $N_u^+$ is in $C(u)$, and there are at least $\Delta_u$ negative edges incident to $u$ with both endpoints in $C(u)$.}\\ 
Let $V^{2a}$ be the vertices of $V^2$ in this case. 
Since these negative edges are disagreements, we may charge $|V_u|$ to $y(u)$:
$$\boxed{\sum_{u \in V^{2a}} |V_u| \leq \frac{10}{3} \cdot \sum_{u \in V} \Delta_u \leq \frac{10}{3}\cdot \sum_u y(u) = \frac{10}{3} \cdot \textsf{OPT}.}$$

\noindent \textbf{Case 2b:} \emph{At least 0.85 fraction of $N_u^+$ is in $C(u)$, and $|C(u)| \leq 2 \cdot \Delta_u$.}\\
Observe Cases 2a and 2b cover the full range of Case 2. 
This is because we can write $|C(u)| = |C(u)\cap N_u^+| + |C(u)\cap N_u^-|$, 
and if there are at \emph{most} $\Delta_u$ negative edges incident to $u$ with both endpoints in $C(u)$, this is equivalent to saying $|C(u)\cap N_u^-| \leq \Delta_u$.
So, both $|C(u)\cap N_u^+|$ and $|C(u)\cap N_u^-|$ are upper bounded by $\Delta_u$ (since $\Delta_u = |N_u^+|$).

\smallskip

Let $V^{2b}$ be the vertices of $V^2$ in this case. 
Let $N_{u,v} := N_v^+ \cap N_u^+ \cap C(u)$. We know from Proposition \ref{prop: pos_overlap} (for $v \in R_1(u)'$ and using the assumption bringing us into Case 2) that $|N_{u,v}| \geq 0.15 \cdot \Delta_u$.  
Observe that for $v \in V_u \subseteq N_u^- \cap \overline{C(u)}$ 
and $w \in N_{u,v}$, the positive edge $(w,v)$ is a disagreement.

Fix $w \in N_u^+ \cap C(u)$ and $u \in V^{2b}$. Define $\varphi(u,w) = |V_u \cap N_w^+|$. We say that $w$ dispenses $\varphi(u,w)/|C(u)|$ charge to $u$. We would like to show that each $u \in V^{2b}$ receives $\Omega(|V_u|)$ charge, and that each $w$ dispenses at most $y(w)$ charge. By double counting,  
$$\sum_{w \in N_u^+ \cap C(u)} \varphi(u,w) \geq |V_u| \cdot 0.15\Delta_u $$
where we have used that $\varphi(u,w)$ is equal to the number of $v \in V_u$ such that $w \in N_{u,v}$, and we know that $|N_{u,v}|$ is at least $0.15 \Delta_u$. So, the total amount of charge that $u \in V^{2b}$ receives is 
$$\sum_{w \in N_u^+ \cap C(u)} \frac{\varphi(u,w)}{|C(u)|} \geq |V_u| \cdot 0.15 \Delta_u / |C(u)| \geq 3/40 \cdot |V_u|$$
since $|C(u)| \leq 2\Delta_u$ by assumption of the case. 

Now we need to upper bound the amount of charge that each $w \in V$ dispenses in total. Note that $\varphi(u,w) \leq y(w)$, and $w$ only dispenses to vertices $u \in C(w) \cap N_w^+$. Further, $C(u) = C(w)$. So $w$ dispenses at most $|C(w)| \cdot y(w)/|C(u)| = y(w)$ charge in total. Putting everything together: 

\begin{align*}
     \sum_{w \in V} y(w) &\geq \text{total charge dispensed} \\
     &\geq \sum_{u \in V^{2b}}  \sum_{w \in N_u^+ \cap C(u)} \frac{\varphi(u,w)}{|C(u)|} 
     \geq \sum_{u \in V^{2b}} \frac{3}{40} \cdot |V_u|.
 \end{align*}
 So in all,
 \[\boxed{\sum_{u \in V^{2b}} |V_u| \leq \frac{40}{3} \cdot \sum_{w \in V} y(w) = \frac{40}{3} \cdot \textsf{OPT}.}\]
This concludes the bounding of the fractional cost of the negative edges, 
and our constant in the approximation factor is $1 + 200/9 + 10/3 + 40/3 < 40$ (adding together the boxed upper bounds). 



\end{proof}

\section{KMZ Rounding Algorithm and Run-time}
\label{sec: lp_rounding_alg}

We construct a fractional solution with our adjusted correlation metric (recall we use this in place of an LP solution), and then
input this solution to a rounding algorithm by Kalhan, Makarychev, and Zhou \citeyearpar{KMZ19}.

First, we set up some notation for the rounding algorithm, 
mainly following the example set by Kalhan, Makarychev, and Zhou.
Let the ball of radius $\rho$, with respect to the distance defined by 
any semi-metric $z$ on $V$, around vertex $u \in V$ 
be denoted $\textsf{Ball}(u,\rho) = \{v \in V \mid z_{uv} \leq \rho\}$.
The algorithm is iterative, and continues to find new clusters until every vertex is clustered.
Let the set of unclustered vertices at time $t$ be denoted by $V_t \subseteq V$.
The cluster center in $V_t$ is the vertex maximizing the quantity
$$
L_t(u)= \sum_{v \in \textsf{Ball}(u,r) \cap V_t} (r-z_{uv})
$$
for $r$ a parameter input to the algorithm.
Note that when $L_t$ is large,
vertices in $\textsf{Ball}(u,r) \cap V_t$ are closely clustered together.\\



\begin{algorithm} \label{KMZ-alg}[Rounding algorithm]
\end{algorithm}

\begin{minipage}{13cm}
\rule{12cm}{0.4pt}\\
\noindent \textbf{Input: }Semi-metric $z$ on $V$.

\noindent \textbf{Output: }Clustering $\mathcal{C}$.
\begin{enumerate}
    \item Let $V_0 = V$, $r = 1/5$, $t=0$.
    \item \textbf{while} ($V_t \neq \emptyset$)
    \begin{itemize}
        \item Find $u^*_t = \arg \max_{u \in V_t} L_t(u) =  \arg \max_{u \in V_t} \sum_{v \in \text{Ball}(u,r) \cap V_t} {r-z_{uv}}$. 
        \item Create ${C_t = \text{Ball}(u^*_t,2r) \cap V_t}$.
        \item Set $V_{t+1} = V_t \setminus C_t$ and $t=t+1$.
    \end{itemize}
    \item Return ${\mathcal{C} = (C_0, \dots, C_{t-1})}$.
\end{enumerate}
    \rule{12cm}{0.4pt}
\end{minipage}\\

We define $\textsf{LP}(u,v)$ to be $(u,v)$'s cost to LP \ref{KMZ_LP} in the objective, i.e., $x_{uv}=\textsf{LP}(u,v)$ for $(u,v) \in E^+$ and 
$x_{uv}=1-\textsf{LP}(u,v)$ if $(u,v) \in E^-$. 
Similarly, we define a cost for each edge with respect to 
Algorithm \ref{KMZ-alg},
$\textsf{ALG}(u,v) = \mathds{1}((u,v) \text { is a disagreement})$.
Then summing over all edges incident to vertex $u$, we obtain the disagreements to each vertex, $\textsf{ALG}(u) = \sum_{v \in V} \textsf{ALG}(u,v)$.

The technical work of Kalhan, Makarychev, and Zhou \citeyearpar{KMZ19} is in showing that 
\begin{equation}\label{eq: KMZ}
    \textsf{ALG}(u) = \sum_{v \in V}\textsf{ALG}(u,v) \leq 
5 \cdot \sum_{v \in V}\textsf{LP}(u,v) = 5 y(u).
\end{equation} 


Davies, Moseley, and Newman \citeyearpar{DMN23} discuss the run-time of Algorithm \ref{KMZ-alg} (see Appendix A in their paper).
The rounding algorithm has run-time $O(n^2)$ for general graphs, 
and the full run-time of the KMZ algorithm is dominated by the time to solve the LP.



